\documentclass[12pt]{article}

\usepackage[margin=1in]{geometry}
\usepackage{amsmath}

\title{Quiz 4: Elect 562}
\author{Jacob Vaught}
\date{}

\begin{document}

\maketitle

\section*{Question 1}
\textbf{Free-Space Path Loss and Different Carrier Frequencies (30 pts)}

The free-space path loss (FSPL) model is used to predict the loss of signal strength of an electromagnetic wave when it travels through free space. FSPL is a function of the distance between the transmitter and receiver and the frequency of the carrier wave. Mathematically, it is expressed as:
\begin{equation}
    \text{FSPL (dB)} = 20\log_{10}(d) + 20\log_{10}(f) + 20\log_{10}\left(\frac{4\pi}{c}\right)
\end{equation}
where \( d \) is the distance between the transmitter and receiver, \( f \) is the frequency of the signal, and \( c \) is the speed of light in a vacuum.

As the equation shows, path loss increases with both an increase in frequency and distance. This occurs because higher frequencies have shorter wavelengths, which leads to greater energy dispersion in space. This dispersion causes the signal's power density to decrease more rapidly as it propagates through space, leading to higher path loss at higher frequencies.

\section*{Question 2}
\textbf{Doubling Carrier Frequency and Coverage (35 pts)}

If an operator doubles the carrier frequency while keeping the antenna gains the same, the coverage is expected to worsen. Here's why:

\textbf{Increased Path Loss (25 pts):} As previously mentioned, higher frequencies experience greater path loss. This means that when the frequency is doubled, the signal will attenuate more rapidly, reducing the effective coverage area.

\textbf{Reduced Wavelength (10 pts):} Doubling the frequency halves the wavelength. Antennas are typically designed for specific wavelengths, so a change in frequency can make the existing antenna less efficient for the new frequency, which could further reduce coverage.

\section*{Question 3}
\textbf{Choosing Carrier Frequency for Robustness Against Mobility (35 pts)}

For a system that needs to be robust against mobility, choosing between 600 MHz (f1) and 5.8 GHz (f2) depends on several factors:

\textbf{Choose 600 MHz (10 pts):} Lower frequencies, like 600 MHz, are generally better for mobility scenarios.

\textbf{Reasons (25 pts):}
\begin{itemize}
    \item \textbf{Better Penetration and Range:} Lower frequencies can penetrate obstacles like buildings more effectively and have a longer range, which is crucial for a mobile environment.
    \item \textbf{Less Path Loss:} As discussed earlier, lower frequencies experience less path loss compared to higher frequencies, leading to more stable signal strength over distances.
    \item \textbf{Less Impact from Doppler Shift:} Mobility introduces Doppler shift, which is less pronounced at lower frequencies, making 600 MHz more stable in high mobility scenarios.
    \item \textbf{Wider Coverage:} With better penetration and range, 600 MHz provides wider coverage, reducing the need for a dense network of base stations, which is advantageous in a mobile setting.
\end{itemize}

In summary, for robustness against mobility, 600 MHz is generally a better choice due to its superior penetration, range, and stability in a mobile environment.

\end{document}
